\chapter{Kidneys}

\label{ch:ch49_kidneys}

\section{What do the kidneys do?}
Your two kidneys sit on either side of your spine, just below the ribcage. They filter waste products from the blood to produce urine, regulate fluid and salt balance, help control blood pressure, and produce hormones that stimulate red blood cell production. Because the kidneys are highly vascular and have a complex internal structure, imaging plays a central role in evaluating a wide range of conditions---from stones and infections to tumors.

\section{When is kidney imaging ordered?}
Doctors order kidney imaging for many reasons:
\begin{itemize}
    \item \textbf{Flank pain or suspected kidney stones:} A non-contrast CT (CT KUB) is the most sensitive test.
    \item \textbf{Unexplained blood in the urine (haematuria):} To look for masses, stones, or structural causes.
    \item \textbf{A lump felt on examination or an abnormality on ultrasound:} To characterize a renal mass.
    \item \textbf{Suspected kidney infection:} Ultrasound to assess for obstruction; CT if complicated pyelonephritis or abscess is suspected.
    \item \textbf{High blood pressure in a young patient:} To exclude renal artery stenosis or parenchymal disease.
    \item \textbf{Known cancer:} Staging and follow-up.
\end{itemize}

\section{What happens during the scan?}
\begin{itemize}
    \item \textbf{Ultrasound:} The kidneys are well visualised because of their posterior location. A full bladder helps assess the urinary tract. No radiation is used.
    \item \textbf{CT KUB (non-contrast):} You lie on a table that passes through the CT scanner. No contrast is needed for stone detection. The scan takes 5--10 minutes.
    \item \textbf{CT urography (contrast):} Contrast is injected and images are acquired in multiple phases (corticomedullary, nephrographic, excretory). This evaluates masses, haematuria, and the urothelium. Takes 15--20 minutes.
    \item \textbf{MRI:} Used for problem-solving, especially for characterising indeterminate masses or when contrast allergy precludes CT. Diffusion-weighted imaging (DWI) improves lesion detection.
\end{itemize}

\section{Normal kidney appearance}
The normal kidney has a smooth contour, distinct corticomedullary differentiation on ultrasound, and homogeneous parenchymal enhancement on CT. The renal pelvis and ureters are seen on excretory-phase imaging. Kidney size ranges from 9--12~cm in adults, with the left kidney slightly larger than the right.

\section{Renal masses: characterisation approach}
When a renal mass is identified, the key question is: \emph{is this a solid enhancing mass (potentially malignant) or a benign entity?}

\begin{table}[h]
\centering
\caption{Common renal masses}
\label{tab:renal_masses}
\begin{tabular}{@{}p{0.20\textwidth}p{0.24\textwidth}p{0.17\textwidth}p{0.22\textwidth}@{}}
\hline
\textbf{Mass} & \textbf{Imaging features} & \textbf{Enhancement} & \textbf{Management} \\
\hline
Angiomyolipoma (AML) & Fat density on CT ($-$10~HU or less) & Variable & Surveillance if $<$4~cm \\
Chromophobe RCC & Large, well-defined & Moderate enhancement & Better prognosis \\
Clear cell RCC & Solid, heterogeneous & Intense, heterogeneous & Staging + nephrectomy \\
Cystic mass & See Bosniak classification & Variable & Bosniak-guided management \\
Oncocytoma & Solid, central scar & Intense, homogeneous & Often indistinguishable from RCC \\
Papillary RCC & Solid, hypovascular & Mild enhancement & Staging + nephrectomy \\
\hline
\end{tabular}
\end{table}

\begin{pearlbox}[Fat is the key]
An adult renal mass containing macroscopic fat is almost always an angiomyolipoma. Confirm by measuring CT attenuation values: fat measures $-30$ to $-100$~HU. Exception: rare fat-containing renal cell carcinomas exist, but they are exceedingly uncommon.
\end{pearlbox}

\section{Bosniak classification of cystic renal masses}
The Bosniak classification (2019 update) stratifies cystic masses by risk of malignancy:

\begin{table}[h]
\centering
\caption{Bosniak classification (v2019)}
\label{tab:bosniak}
\begin{tabular}{@{}p{0.14\textwidth}p{0.39\textwidth}p{0.14\textwidth}p{0.17\textwidth}@{}}
\hline
\textbf{Category} & \textbf{Description} & \textbf{Malignancy risk} & \textbf{Management} \\
\hline
I & Simple cyst, hairline wall, no septa/nodules & 0\% & No follow-up \\
II & Minimally complex: $\leq$3 thin septa, $\leq$3~cm hyperdense nodule & $<$1\% & No follow-up \\
IIF & $\geq$4 thin septa or thickened smooth septa/nodule, $\geq$3~cm & 5--10\% & Imaging surveillance \\
III & Thickened irregular septa or enhancing nodules & 50--55\% & Surgery or ablation \\
IV & Clearly malignant features: solid enhancing components & $>$90\% & Surgery \\
\hline
\end{tabular}
\end{table}

\section{Renal mass protocol CT}
A dedicated renal mass protocol includes:
\begin{enumerate}
    \item \textbf{Pre-contrast:} Establishes baseline attenuation and detects fat or hemorrhage.
    \item \textbf{Corticomedullary phase (30--40~s):} Optimises detection of hypervascular lesions (clear cell RCC) and delineates renal vasculature for surgical planning.
    \item \textbf{Nephrographic phase (80--120~s):} Most sensitive phase for solid mass detection; optimises corticomedullary differentiation for staging.
    \item \textbf{Excretory phase (5--10~min):} Evaluates the collecting system and urothelium for transitional cell carcinoma.
\end{enumerate}

\begin{pearlbox}[Enhancement = malignancy until proven otherwise]
A renal mass that enhances by $>$15~HU after contrast (i.e., absolute increase of $>$15~HU from the pre-contrast phase) should be considered a solid enhancing mass and managed accordingly. Enhancement of less than 10~HU is generally considered non-enhancing.
\end{pearlbox}

\section{Partial vs radical nephrectomy: imaging role}
Imaging helps the urologist decide between nephron-sparing (partial) and radical nephrectomy:
\begin{itemize}
    \item \textbf{Favour partial nephrectomy:} Tumor $<$4~cm (T1a), exophytic location, away from hilum, solitary kidney, bilateral tumors, or underlying renal insufficiency.
    \item \textbf{Favour radical nephrectomy:} Tumor $>$7~cm (T2), central/hilar location, invasion of renal vein or collecting system, or multifocal disease.
    \item \textbf{Imaging staging} includes assessment of renal vein tumor thrombus (coronal/sagittal reformats), perinephric fat stranding, lymphadenopathy, and distant metastases (lungs on CT chest).
\end{itemize}

\section{Hydronephrosis}
Hydronephrosis is dilation of the renal collecting system due to urinary obstruction. Ultrasound is the first-line investigation.

\begin{table}[h]
\centering
\caption{Grading of hydronephrosis}
\label{tab:hydronephrosis}
\begin{tabular}{lll}
\hline
\textbf{Grade} & \textbf{Description} & \textbf{Significance} \\
\hline
0 & Indented renal pelvis & Normal \\
1 & Mild pelvicalyceal dilation & May be physiological \\
2 & Dilation of pelvis and calyces & Likely significant obstruction \\
3 & Rounded calyces, thinned parenchyma & Significant obstruction \\
4 & Gross dilation, cortical thinning & Chronic, severe obstruction \\
\hline
\end{tabular}
\end{table}

Common causes of obstruction include ureteric stones, extrinsic compression (retroperitoneal fibrosis, pelvic malignancy), ureteric strictures, and pelvi-ureteric junction (PUJ) obstruction.

\section{Renal stones}
Most renal stones are radiopaque and well seen on non-contrast CT. Stones appear as high-attenuation foci ($>$400~HU) within the collecting system. Uric acid stones are less dense ($<$500~HU) and may be visible on ultrasound but not on plain radiography. The CT KUB protocol involves a limited scan from the kidneys to the symphysis pubis, reducing radiation exposure compared to a full abdomen--pelvis CT.

\begin{pearlbox}[Size matters for management]
Stones $<$5~mm have a $>$90\% chance of spontaneous passage. Stones 5--10~mm may pass but often require intervention. Stones $>$10~mm or those causing complete obstruction, intractable pain, or infection typically require urological intervention (ureteroscopy, PCNL, or ESWL).
\end{pearlbox}

\section{Renal infection}
Pyelonephritis typically shows perinephric fat stranding, renal enlargement, and heterogeneous enhancement on CT. Focal areas of non-enhancement may mimic a mass (focal pyelonephritis). Renal abscess appears as a rim-enhancing fluid collection. Emphysematous pyelonephritis---a gas-forming infection---is a urological emergency requiring urgent drainage and antibiotics.

\section{Renal vein thrombosis}
Renal vein thrombosis (RVT) is recognised as a filling defect within the renal vein on contrast-enhanced CT or MRI. Causes include nephrotic syndrome, malignancy, dehydration, hypercoagulable states, and post-transplant. Imaging may also show renal enlargement, perinephric edema, and collateral venous pathways in chronic cases.

\section{Reporting checklist}
A structured renal report should include:
\begin{enumerate}
    \item Kidney size, position, and symmetry
    \item Parenchymal thickness and enhancement
    \item Collecting system dilation (grade)
    \item Focal masses: size, location (Bosniak or solid characterisation), enhancement
    \item Stones: location, size, degree of obstruction
    \item Ureter: course, dilatation, filling defects
    \item Renal vasculature (artery and vein patency)
    \item Perinephric and retroperitoneal assessment
    \item Adrenal glands
    \item Comparison with prior imaging
\end{enumerate}

\clearpage
